% Käsitsi ümber kirjutatud kogumiku lahendusest L20 (lk 40).
%
% KAKS PARANDUST ORIGINAALI SUHTES:
% 1. Kogumikus on kirjas "alpha = pi/sqrt(g)", mis annaks 1,00 s/m^0,5.
%    Õige on 2*pi/sqrt(g) = 2,007 s/m^0,5 ja kogumik ise ütleb samas lauses
%    "~ 2 s/m^0,5", seega tegu on trükiveaga. Parandatud.
% 2. Hindamisskeemis on "alpha = T iga katseseeria jaoks"; ruutjuur kadus
%    PDF-i tekstikihist. Õige on alpha = T/sqrt(l).
Katse läbiviimiseks tuleb niidijupp siduda mutri külge ja niidi teine ots
fikseerida kinnitusklambri vahele. Seejärel tuleb mõõta niidi pikkus (mutri
massikeskmest kinnituskohani) ja mõõta võnkeperiood. Võnkeperioodi täpseks
määramiseks tuleb stopperiga mõõta aeg, mis kulub mitme võnke tegemiseks
($> 5$), ja jagada see võngete arvuga.

Konstandi $\alpha$ võimalikult täpseks hindamiseks tuleb katseid korrata.
Niidi pikkust võib, kuid ei pea varieerima. Vastus tuleb matemaatilise
pendli teooriast:
\[
T = 2\pi\sqrt{\frac{l}{g}} = \alpha\sqrt{l}
\quad\Rightarrow\quad
\alpha = \frac{2\pi}{\sqrt{g}} \approx \SI{2,0}{s/m^{0,5}}
\]
Korrektselt läbi viidud katse tulemus ei erine sellest rohkem kui
\SI{10}{\percent}.

\textbf{Hindamine} (kokku 10 p): katse idee \pp{2}; võnkeperioodi mõõtmisel
$n$ võnkeks ($> 5$) kulunud aja lugemine ja sellest võnkeperioodi
arvutamine \pp{1} (kui on mõõdetud ainult üheks võnkeks kulunud aeg, siis
seda punkti ei saa); korduskatsed \pp{$3\times 1$} (iga korduskatse annab ühe
punkti, maksimaalselt 3); arvutus $\alpha = T/\sqrt{l}$ iga katseseeria
jaoks \pp{$3\times 0{,}5$}, mitme mõõtmise korral $\alpha$ keskmistamine
\pp{0,5}; realistlik vastus ($\pm\SI{10}{\percent}$) \pp{2}.
